\section{Method}
\label{sec:method}

\subsection{A Small Model}

We use a deliberately simple three-stage computation. Let $x \in
\mathbb{R}^{d}$ be an input vector and let $W$ be a learned projection. An
aligned environment is a useful stress case because it contains several
visual rows but should remain one coherent editable formula:
\begin{align}
  z &= Wx + b, \\
  \hat{y} &= \operatorname{softmax}(z), \\
  \mathcal{L} &= -\sum_{k=1}^{K} y_k \log \hat{y}_k.
\end{align}

The following equation has a label and a reference, which lets the browser be
tested with both rendered mathematics and cross-references:
\begin{equation}
  \operatorname{score}(x) = \frac{1}{K}\sum_{k=1}^{K}\hat{y}_k.
  \label{eq:score}
\end{equation}

Equation~\ref{eq:score} is not intended as a scientific claim. It is a stable
anchor for testing how a rendered block behaves when a human selects nearby
prose or opens the formula editor.

\subsection{Implementation Notes}

This fixture intentionally uses a multi-file source layout. The renderer must
follow the existing \texttt{\\input} graph while retaining each included file
as the source-location hint. PairTeX does not rewrite the graph or introduce a
second document format.
